%Predložak za seminarski rad%
%Glavna fajla%
%
\documentclass[a4paper, oneside, 12pt]{book}
\usepackage[utf8]{inputenc} %prepoznavanje dijakritičkih znakova
\usepackage[T1]{fontenc} %unos dijakritičkih znakova izravno s tipkovnice
\usepackage[T2A]{fontenc} %унос ћирилићних знакова директно са тастатуре
\usepackage[russian, english, croatian]{babel} %ruski je važan radi nekim matematskih funkcija
\addto\captionscroatian{\renewcommand{\tablename}{Tabela}}
\usepackage{indentfirst} %uvlačenje prvog paragrafa u odjeljku
\usepackage{titlesec} %definiranje proreda između raznih naslova
\usepackage{amsmath, amssymb, amsfonts, esint} %osnovni matematički mod ako se koristi lmodern
\usepackage{mathrsfs} %matematički mod s klasičnim pisanim slovima, \mathscr{}
\usepackage[bookmarks, bookmarksnumbered, hidelinks, unicode, pdfpagelayout=SinglePage]{hyperref} %bookmark pdf
\usepackage{siunitx} %međunarodni sustav mjernih jedinica
\usepackage[a4paper, left=25mm, right=25mm, top=25mm, bottom=25mm]{geometry} %dimenzije papira
\usepackage{graphicx} %napredni grafički mod
\usepackage{color} %definiranje boje simbola i pozadine
\usepackage[justification=centering]{caption} %centriranje opisa slika i tablica
\usepackage{float} %podešavanje pozicija slika i tablica
\usepackage[shortlabels]{enumitem} %podešavanje nabrajanja (slova, brojevi, ...)
\usepackage{array} %tablice
\usepackage{setspace} %podešavanje razmaka teskta
\usepackage{icomma} %zarez kao decimalna razdvojnica
\usepackage[square,sort,comma,numbers]{natbib} %podešavanje razmaka bibliografije
\usepackage{fancyhdr} %manipuliranje zaglavljem i podnožjem
\usepackage{tikz}
\usetikzlibrary{arrows.meta,positioning}
\fancypagestyle{plain} %potrebno zbog numeriranja prve stranice
\fancyhead{} %brisanje zaglavlja
\fancyfoot{} %brisanje podnožja
\cfoot{\footnotesize{\thepage}} %numeriranje prve stranice
\pagestyle{fancy} %uvijek u kombinaciji s fancyhdr
\fancyhead{} %brisanje zaglavlja
\fancyfoot{} %brisanje podnožja
\fancyhead[C]{\scriptsize{SEMINARSKI RAD: Adi Kadušić}} %zaglavlje
\fancyfoot[C]{\footnotesize{\thepage}} %podnožje
\renewcommand{\headrulewidth}{0.4pt} %širina linije za odvajanje zaglavlja
\renewcommand{\footrulewidth}{0pt} %širina linije za odvajanje podnožja
\newcommand{\bunderline}[1]{\underline{#1\mkern-4mu}\mkern4mu} %podešava dužinu donje crte
\newcommand{\boverline}[1]{\overline{#1\mkern-4mu}\mkern4mu} %podešava dužinu gornje crte
\newlist{inlinelist}{enumerate*}{1} %postavljanje "inline" liste
\setlist*[inlinelist,1]{label=\textit{\roman*}),} %podešavanje stila "inline" liste
\setlength{\headheight}{10 mm} %širina zaglavlja
\setlength{\parindent}{0.5 cm} %dubina uvlačenja prve riječi paragrafa
\setlength{\parskip}{0.5 em} %veličina proreda između paragrafa
\setlength{\bibsep}{6 pt} %veličina proreda između bibliografksih referenci
\titleformat{\chapter}{\normalfont\huge\bf}{\thechapter.}{12pt}{\huge\bf} %naslov poglavlja ...
\titleformat{\section}{\normalfont\Large\bf}{\thesection.}{12pt}{\Large\bf} %odjeljka ...
\titleformat{\subsection}{\normalfont\large\bf}{\thesubsection.}{12pt}{\large\bf} %pododjeljka
\titlespacing\section{0pt}{3em}{1em} %prored oko naslova odjeljka
\titlespacing\subsection{0pt}{2em}{1em} %prored oko naslova pododjeljka

%
\begin{document}
%
\frontmatter
	% !TEX root = Glavna.tex
%Predložak za seminarski rad%
%Naslovna stranica%
%
%
\begin{titlepage}
	\begin{center}
		\normalsize{UNIVERZITET U ZENICI}\\
		\textbf{\large{POLITEHNIČKI FAKULTET}}\\
		\textbf{{Studijski program: SOFTVERSKO INŽENJERSTVO}}\\ [9cm]
		\textbf{\huge{Bit-banged UART\\\vspace{0.5ex}u AVR asemblerskom jeziku}}\\
		\textsc{\footnotesize{Automati i formalni jezici}}\\ [9cm]
		\noindent
		\begin{minipage}{0.45\textwidth}
			\begin{flushleft}
				\hspace{1.5 cm}
				\normalsize{Nastavnik:}\\
				red.\,prof.\,dr.\,sc.\,Alen Begović
			\end{flushleft}
		\end{minipage}
		\begin{minipage}{0.45\textwidth}
			\begin{flushright}
				\hspace{2cm}
				\normalsize{Student:\hspace{8mm}\textcolor{white}{.}}\\
			 	Adi Kadušić,\,II-134
			\end{flushright}
		\end{minipage}
		\vfill
		\normalsize{Zenica,\space\the\year.\,godina}
	\end{center}
\end{titlepage}
\newpage
	\phantomsection
	\addcontentsline{toc}{chapter}{Sažetak}
	%Predložak za seminarski rad%
%Naslovna stranica%
%
%
\vspace*{5 em}

\normalsize\noindent\emph{Sažetak}: U ovom radu opisuje se implementacija
asinhronog serijskog \foreignlanguage{english}{UART} protokola na
mikrokontroleru \foreignlanguage{english}{ATmega328P}, izvedena
isključivo softverski, bez upotrebe ugrađenog hardverskog
\foreignlanguage{english}{UART} perifernog modula, pristupom poznatim
kao \foreignlanguage{english}{bit-banging}. Aktuelnost teme obrazlaže se
činjenicom da se \foreignlanguage{english}{UART} protokol, uprkos svojoj
starosti, i dalje široko koristi za komunikaciju sa mikrokontrolerima i
perifernim uređajima, dok se njegov format okvira može posmatrati kao
pogodan primjer konačnog automata čije je stanje potrebno precizno
vremenski uskladiti sa spoljnim signalom, bez oslonca na zajednički
takt predajnika i prijemnika. Cilj rada je da se, na osnovu poznavanja
broja ciklusa izvršavanja pojedinačnih instrukcija
\foreignlanguage{english}{AVR} arhitekture, pokaže na koji način se
konstruiše potprogram za precizno kašnjenje kojim se ostvaruje zadata
brzina prijenosa, te da se na osnovu njega implementiraju potprogrami za
slanje i prijem pojedinačnog bajta podataka. Pri izradi rada korištena
je analiza opkoda i broja ciklusa korištenih instrukcija, izvođenje
matematičke formule za broj ponavljanja petlje kašnjenja, kao i
funkcionalno testiranje realizovanog programa na fizičkoj razvojnoj
ploči.

\noindent\small{\textbf{\emph{Ključne riječi}: ATmega328P, AVR asembler, bit-banging, kašnjenje, mikrokontroler, UART protokol}}.

\vspace{5 em}

\normalsize\noindent\foreignlanguage{english}{\emph{Abstract}: In this paper, the implementation of the asynchronous serial UART protocol on the ATmega328P microcontroller is described, carried out entirely in software, without the use of the built-in hardware UART peripheral module, an approach commonly referred to as bit-banging. The relevance of the topic is explained by the fact that the UART protocol, despite its age, is still widely used for communication with microcontrollers and peripheral devices, while its frame format can be regarded as a suitable example of a finite state machine whose state must be precisely timed against an external signal, without relying on a shared clock between the transmitter and the receiver. The aim of the paper is to show, based on knowledge of the number of execution cycles of individual AVR instructions, how a subroutine for precise delay is constructed to achieve the required baud rate, and how it is subsequently used to implement the transmission and reception of a single data byte. The work is based on an analysis of instruction opcodes and cycle counts, the derivation of a mathematical formula for the number of delay loop iterations, and functional testing of the resulting program on a physical development board.}

\noindent\small{\foreignlanguage{english}{\textbf{\emph{Keywords}}: ATmega328P, AVR assembler, baud rate, bit-banging, delay, microcontroller, UART protocol}}.
\normalsize
\newpage
	\phantomsection
	\begin{spacing}{0.5}
		\tableofcontents
	\end{spacing}
	\addcontentsline{toc}{chapter}{Sadržaj}	
	% Uvod
%
%
\phantomsection
\addcontentsline{toc}{chapter}{Uvod}
\chapter*{Uvod}
\label{chap:intro}
%
%
Serijska komunikacija smatra se jednim od najstarijih i
najrasprostranjenijih načina razmjene podataka između mikrokontrolera i
perifernih uređaja, a \foreignlanguage{english}{UART}
(\foreignlanguage{english}{Universal Asynchronous Receiver-Transmitter})
protokol je zbog svoje jednostavnosti i malog broja potrebnih linija
ostao u širokoj upotrebi i danas, uprkos postojanju brojnih novijih
protokola. Posebno zanimljivim svojstvom ovog protokola smatra se
njegova asinhronost, odnosno odsustvo zajedničke linije takta između
predajnika i prijemnika, zbog čega se tačan trenutak očitavanja svakog
pojedinačnog bita mora unaprijed izračunati na osnovu poznate brzine
prijenosa. Format okvira, koji se sastoji od tačno određenog niza stanja
kroz koja linija prolazi u tačno određenom vremenskom rasporedu, smatra
se pogodnim praktičnim primjerom konačnog automata čije prelaze između
stanja ne pokreće samo promjena ulaznog signala, već i protok vremena.

Iako su savremeni mikrokontroleri, uključujući i
\foreignlanguage{english}{ATmega328P} korišten u ovom radu, opremljeni
ugrađenim hardverskim \foreignlanguage{english}{UART} perifernim
modulom kojim se ova komunikacija izvodi automatski, kod softverske
realizacije istog protokola, poznate kao \foreignlanguage{english}{bit-banging},
tajming svake pojedinačne operacije mora se precizno kontrolisati na
nivou broja ciklusa procesora. Zbog ove osobine, softverska realizacija
smatra se pogodnom za razumijevanje veze između arhitekture procesora,
skupa njegovih instrukcija i vremenski osjetljivih protokola opisanih
formalnim automatima. Svrha ovog rada je da se na konkretnom primjeru
pokaže postupak kojim se, polazeći od poznatog broja ciklusa svake
korištene instrukcije, konstruiše program koji ostvaruje zadatu brzinu
prijenosa isključivo korištenjem digitalnih ulazno-izlaznih pinova.

U radu će se odgovoriti na pitanje kako se, bez upotrebe hardverskog
\foreignlanguage{english}{UART} modula, može ostvariti tačno određena
brzina prijenosa korištenjem petlje za kašnjenje čiji je broj ponavljanja
izveden iz broja ciklusa procesora. Razmotrit će se koje su instrukcije
\foreignlanguage{english}{AVR} arhitekture najpogodnije za rad sa
pojedinačnim bitom ulazno-izlaznog registra i na koji način njihov izbor
utiče na preciznost i predvidljivost tajminga. Poseban naglasak bit će
stavljen na pitanje kako se ista petlja za kašnjenje, uz odgovarajuće
prilagođavanje okolnog koda, može iskoristiti i za slanje i za prijem
podataka, uprkos tome što se struktura koda oko poziva te petlje
razlikuje u oba slučaja.

Rad se sastoji iz tri poglavlja. U prvom poglavlju izlaže se teorijska
osnova neophodna za razumijevanje ostatka rada. Opisuje se format
\foreignlanguage{english}{UART} okvira i način izvođenja brzine
prijenosa, kao i razlog očitavanja na sredini bita, a zatim se uvode
dijelovi \foreignlanguage{english}{AVR} arhitekture koji se koriste u
implementaciji, uključujući registre opće namjene, status registar
\foreignlanguage{english}{SREG}, ulazno-izlazne portove i trajanje
izvršavanja instrukcija u ciklusima.

U drugom poglavlju izvodi se vremenska analiza programa na nivou
pojedinačnih instrukcija. Opisuje se skup korištenih instrukcija zajedno
sa njihovim opkodom i tačnim brojem ciklusa izvršavanja, a zatim se, na
osnovu tako utvrđenih vrijednosti, izvodi opšta formula kojom se
određuje broj ponavljanja petlje za kašnjenje. Formula se potom
primjenjuje na dva konkretna slučaja, punu bit periodu i poravnanje na
sredinu bita prilikom prijema.

U trećem poglavlju opisuje se praktična implementacija programa.
Objašnjava se sintaksa i skup direktiva korištenih pri pisanju izvornog
koda, zatim algoritam slanja i prijema podataka uz obrazloženje odabira
korištenih instrukcija i pomoćnih makroa za usklađivanje broja ciklusa.
Na kraju poglavlja opisuje se razvojno okruženje i postupak testiranja
realizovanog programa na fizičkoj razvojnoj ploči.

Na kraju rada nalazi se zaključak, u kojem su sistematizovani
najvažniji rezultati i doprinosi rada kao cjeline. Pored zaključka, u
radu se nalaze i popis slika, kao i lista korištenih referenci.

\newpage

%
\mainmatter
	%Poglavlje 1
%
%
\chapter{Teorijska osnova}
\label{chap:theory}
%
U ovom poglavlju izlaže se teorijska osnova neophodna za razumijevanje
implementacije opisane u narednim poglavljima. Prvo se objašnjava format
UART okvira i način na koji se iz njega izvodi brzina prijenosa podataka,
kao i razlog zbog kojeg se prijem vrši na sredini bita. Zatim se opisuju
dijelovi AVR arhitekture koji se koriste u implementaciji: registri opće
namjene, status registar SREG, ulazno-izlazni portovi te trajanje
izvršavanja instrukcija u taktovima.
%
\section{UART protokol}
\label{sec:uart}
%
UART je asinhroni serijski protokol, što znači da predajnik i prijemnik ne
dijele zajedničku liniju takta. Podaci se prenose u okvirima koji se
sastoje od start bita, osam podatkovnih bita i stop bita. Linija je u
mirovanju (\foreignlanguage{english}{idle}) na visokom naponskom nivou, a
start bit se signalizira prelaskom linije na nizak nivo, čime prijemnik
prepoznaje početak novog okvira. Nakon start bita slijedi osam podatkovnih
bita, pri čemu se prvo prenosi bit najmanje težine
(\foreignlanguage{english}{LSB}). Okvir se završava stop bitom, tokom
kojeg se linija vraća na visok nivo, čime se osigurava razmak prije
eventualnog sljedećeg okvira.

Brzina prijenosa (\foreignlanguage{english}{baud rate}) predstavlja broj
bita koji se prenese u jednoj sekundi i definiše se kao recipročna
vrijednost trajanja jednog bita, odnosno $B = 1/T_{bit}$, gdje je
$T_{bit}$ trajanje jednog bita izraženo u sekundama. Pošto predajnik i
prijemnik ne dijele zajednički takt, oba uređaja moraju unaprijed
poznavati istu vrijednost $T_{bit}$ kako bi tajming predaje i prijema bio
usklađen.

Pošto UART protokol ne koristi zajedničku liniju takta, prijemnik mora
sam odrediti trenutak u kojem će očitati vrijednost svakog bita, na
osnovu poznate brzine prijenosa i trenutka detekcije start bita. Ivice
bita, odnosno prelazi između susjednih bita, najosjetljivije su na šum i
na malo odstupanje takta između predajnika i prijemnika, pa bi očitavanje
na samoj ivici moglo dati pogrešnu vrijednost. Očitavanjem na sredini
bita ostavlja se maksimalna vremenska margina sa obje strane, čime se
povećava otpornost prijema na akumulirano odstupanje tajminga tokom
trajanja okvira.
%
\section{AVR arhitektura}
\label{sec:avr}
%
ATmega328P mikrokontroler sadrži 32 8-bitna registra opće namjene, označena od R0 do R31,
organizovana u registar
\foreignlanguage{english}{file} koji je direktno povezan sa
aritmetičko-logičkom jedinicom. Ovakva organizacija omogućava da se
većina aritmetičkih i logičkih operacija izvrši u jednom ciklusu takta,
jer se oba operanda čitaju iz registar \foreignlanguage{english}{file}-a
i rezultat se u istom ciklusu upisuje nazad. Ovi registri su namijenjeni
isključivo za privremeno čuvanje podataka i izvođenje proračuna unutar
CPU jezgra i sami po sebi nemaju nikakav uticaj na fizičke pinove
mikrokontrolera.

Za upravljanje perifernim jedinicama, uključujući i digitalne ulazno-izlazne
pinove, mikrokontroler posjeduje zaseban adresni prostor
ulazno-izlaznih (I/O) registara, fizički odvojen od registar
\foreignlanguage{english}{file}-a registara opće namjene. Dok registri
opće namjene služe isključivo procesoru za privremeno čuvanje podataka,
I/O registri predstavljaju direktno preslikavanje stanja i konfiguracije
hardvera, poput porta, i svaka promjena njihove vrijednosti odmah se
odražava na fizičkom pinu. Ovaj zaseban prostor takođe omogućava da se
pojedinačni bit u I/O registru postavi ili obriše bez potrebe da se
prethodno pročita i ponovo upiše cijeli sadržaj registra, što je posebno
značajno kada je potrebno precizno kontrolisati trajanje operacije u
ciklusima.

Svaki digitalni I/O port, na primjer port D, kontroliše se pomoću tri
odvojena I/O registra. Registar DDRx određuje smjer svakog pojedinačnog
pina tog porta, odnosno da li se pin ponaša kao ulaz ili izlaz. Registar
PORTx ima dvostruku ulogu koja zavisi od stanja odgovarajućeg bita u
DDRx: kada je pin konfigurisan kao izlaz, bit u PORTx direktno određuje
da li je na pinu prisutan nizak ili visok naponski nivo, dok kod
ulaznog pina isti bit uključuje ili isključuje interni pull-up otpornik.
Registar PINx je, za razliku od druga dva, namijenjen isključivo čitanju
i uvijek odražava stvarni trenutni naponski nivo prisutan na fizičkom
pinu, bez obzira na to da li je taj pin konfigurisan kao ulaz ili izlaz.
Na primjer, da bi se neki pin porta D postavio u stanje logičke
jedinice, potrebno je najprije u registru DDRD postaviti odgovarajući
bit na jedan, čime se taj pin konfiguriše kao izlazni, a zatim u
registru PORTD postaviti isti bit na jedan, čime se na tom pinu
uspostavlja visok naponski nivo.

Status registar (SREG) je poseban I/O registar širine 8 bita, u kojem
svaki pojedinačni bit predstavlja zasebnu zastavicu koja nakon izvršenja
većine aritmetičkih i logičkih operacija bilježi neko svojstvo dobijenog
rezultata. Registar sadrži sljedećih osam zastavica:

\begin{itemize}
    \item C (\foreignlanguage{english}{Carry}) - zastavica prijenosa,
        postavlja se kada operacija proizvede prijenos ili posudbu izvan
        granice od 8 bita;
    \item Z (\foreignlanguage{english}{Zero}) - postavlja se kada je
        rezultat posljednje operacije jednak nuli;
    \item N (\foreignlanguage{english}{Negative}) - odražava najznačajniji
        bit rezultata i pokazuje da li je rezultat negativan;
    \item V (\foreignlanguage{english}{Overflow}) - zastavica
        prekoračenja opsega kod aritmetike sa predznakom;
    \item S (\foreignlanguage{english}{Sign}) - stvarni predznak
        rezultata, dobijen kao isključivo ILI zastavica N i V;
    \item H (\foreignlanguage{english}{Half Carry}) - prijenos između
        nižeg i višeg nibla bajta, koristan kod BCD aritmetike;
    \item T (\foreignlanguage{english}{Bit Copy Storage}) - jednobitno
        mjesto za privremeno čuvanje vrijednosti jednog bita prilikom
        prenošenja između registara;
    \item I (\foreignlanguage{english}{Global Interrupt Enable}) -
        globalno omogućava ili onemogućava obradu prekida.
\end{itemize}

Sadržaj ovog registra se može ispitati kako bi program donio odluku o
daljem toku izvršavanja, na primjer da li će se izvršiti uslovni skok,
čime SREG postaje osnovni mehanizam preko kojeg se ostvaruje uslovno
grananje u programu.

Većina AVR instrukcija izvršava se u jednom ciklusu takta, zahvaljujući
jednostepenom pipeline-u u kojem se dohvatanje sljedeće instrukcije
preklapa sa izvršavanjem trenutne. Instrukcije koje mijenjaju tok
programa, poput relativnih skokova i poziva potprograma, instrukcije
koje ispituju bit i uslovno preskaču sljedeću instrukciju, te uslovni
skokovi kada se grananje zaista izvrši, zahtijevaju dodatni ciklus ili
više, jer se u tim slučajevima pipeline mora isprazniti ili se dodatno
pristupa memoriji programa odnosno steku.
%

	%Poglavlje 2
%
%
\chapter{Opcode i cycle analiza instrukcija}
\label{chap:opcode}
%
\section{Instrukcije i njihov tajming}
\label{sec:instructions}
%
abc
%
\section{Primjena: izvod formule za delay}
\label{sec:delay}
%
abc
%
	%Poglavlje 3
%
%
\chapter{Implementacija}
\label{chap:implementation}
%
\section{Sintaksa i direktive AVR asemblera}
\label{sec:syntax}
%
abc
%
\section{Algoritam i kod}
\label{sec:algorithm}
%
abc
%
\section{Testiranje}
\label{sec:testing}
%
abc
%


%
\backmatter
	%Predložak za seminarski rad%
%Zaključak%
%
%
\phantomsection
\addcontentsline{toc}{chapter}{Zaključak}
\chapter*{Zaključak}
\label{chap:conc}
%
U ovom radu realizovan je softverski, odnosno
\foreignlanguage{english}{bit-banged}, \foreignlanguage{english}{UART}
protokol na mikrokontroleru \foreignlanguage{english}{ATmega328P}, bez
upotrebe njegovog ugrađenog hardverskog \foreignlanguage{english}{UART}
modula. Na osnovu opkoda i tačnog broja ciklusa korištenih instrukcija
izvedena je opšta formula za broj ponavljanja petlje za kašnjenje, koja
je zatim primijenjena na punu bit periodu i na poravnanje na sredinu
bita prilikom prijema, čime je teorijski određen broj ponavljanja
petlje potreban za ostvarivanje zadate brzine prijenosa.

Ispravnost implementacije provjerena je funkcionalnim testiranjem na
fizičkoj razvojnoj ploči Arduino Uno R3, uz korištenje alata \emph{avra}
za asembliranje i alata \emph{avrdude} za programiranje mikrokontrolera.
Komunikacijom uspostavljenom preko terminalskog alata \emph{tio}
utvrđeno je da program ispravno prima i vraća nazad svaki poslani
karakter, čime je potvrđena ispravnost implementiranih potprograma za
slanje i prijem podataka u punom krugu.

Testiranje je, međutim, izvedeno isključivo na ovaj funkcionalan način,
bez upotrebe osciloskopa ili logičkog analizatora, tako da stvarno
trajanje pojedinačnog bita na predajnoj liniji nije nezavisno izmjereno.
Kao pravac daljeg rada preporučuje se upravo takva empirijska provjera,
kao i proširenje implementacije mogućnošću podešavanja brzine prijenosa
i detekcijom greške okvira, čime bi se ostvarena rješenja mogla porediti
sa ponašanjem hardverskog \foreignlanguage{english}{UART} modula pod
istim uslovima.
\newpage

	\phantomsection
    \addcontentsline{toc}{chapter}{Popis slika}
	\begin{spacing}{0.4}
		\listoffigures
	\end{spacing}
	\cleardoublepage
	\phantomsection
	\bibliographystyle{unsrtnat}
	\addcontentsline{toc}{chapter}{Bibliografija}
	\nocite{*}
	\bibliography{Reference}
%
\end{document}